\documentclass[10pt,a4paper]{amsart}
\usepackage[margin=19mm]{geometry}
\usepackage{amsmath,amssymb,mathtools}
\usepackage[scr=rsfs,cal=euler]{mathalfa}
\usepackage{xfrac}
\usepackage[unicode,hidelinks,bookmarks=false]{hyperref}
\newcommand{\ie}{i.e.\ }
\newcommand{\cf}{cf.\ }
\newcommand{\resp}{respectively\ }
\newcommand{\Subsection}{\subsection}
\providecommand{\nobreakdash}{\nobreak\hbox{-}}
\setlength{\emergencystretch}{2em}
\multlinegap0pt
\usepackage{csquotes}
\usepackage{amssymb}
\usepackage{xcolor}
\newtheorem{theorem}{Theorem}
\title{Mitchell Rank for Supercompactness}
\author{Erin Carmody}
\date{Source: arXiv:2602.08852v2 (CC BY 4.0)}
\begin{document}
\maketitle

\noindent\emph{Source and licence.} Mitchell Rank for Supercompactness. Version arXiv:2602.08852v2 (CC BY 4.0), distributed by arXiv under the Creative Commons Attribution 4.0 International licence, http://creativecommons.org/licenses/by/4.0/, as recorded in the arXiv licence metadata for this exact version and shown on the abstract page https://arxiv.org/abs/2602.08852v2. Full licence notice: \texttt{licenses/arxiv-2602.08852-CC-BY-4.0.txt}.\par\smallskip

\noindent\textit{Editorial scope.} Counted unit: the article's theorem theorem (source0.tex lines 284--285). The article's complete original proof of it is delivered with it (source0.tex lines 287--299, one \texttt{proof} environment of 13 lines). No mathematics is written, repaired or re-derived: the statement and the proof are the article's own lines, in the article's own notation, and no retained line is substituted or re-flowed.  No further source-local context is needed: every cross-reference of every retained line resolves inside the two delivered spans.  Nothing was omitted from any retained span, so the package carries no omission record.  No retained line contains a citation, so the wrapper carries no bibliography and no bibliography entry.  No retained line contains a figure environment, an image-inclusion command or drawing code, so no image is delivered and the package declares no asset dependency.  Editorial formatting only, in addition: an AMS \texttt{amsart} preamble that restates the article's own declarations and macros used by the retained lines, namely the environments \texttt{theorem} on separate counters, together with the standard packages those lines need.  The retained environments print in this document as Theorem 1 in order of appearance, whereas the article prints the same items with the numbering of its own full document.  The complete native source of the article as distributed in the arXiv e-print for this exact version is bundled unchanged as \texttt{sources/194278/source0.tex}.
\begin{theorem} For any $V \models ZFC + GCH$, any $\Theta: ORD \to ORD$ and any $F: ORD \to ORD$, there is a forcing extension $V[G]$ where, if $\kappa$ is $\Theta(\kappa)$-supercompact, $\Theta \upharpoonright \kappa$ represents $\Theta(\kappa)$, $F \upharpoonright \kappa$ represents $F(\kappa)$ in $V$, and $\Theta''\kappa \subseteq \kappa$, then $o_{\Theta(\kappa)-sc}(\kappa)^{V[G]} = \min\{o_{\Theta(\kappa)-sc}(\kappa)^V, F(\kappa)\}$.
\end{theorem}

\begin{proof} Suppose $V \models$ ZFC + GCH. Let $\Theta:$ ORD $\to$ ORD and $F:$ ORD $\to$ ORD. Let $\mathbb{P}$ be an Easton support ORD-length iteration which forces at stage $\gamma$ to add a club $c_{\gamma} \subseteq \gamma$ such that $\delta \in c_{\gamma}$ implies $o_{\Theta(\delta)-sc}(\delta)^V < F(\delta)$, whenever $\gamma$ is a a closure point of $\Theta$, meaning $\Theta''\gamma \subseteq \gamma$, otherwise force trivially. Let $G \subseteq \mathbb{P}$ be $V$-generic. Let $\delta_0$ be the first inaccessible cardinal. The forcing before stage $\delta_0$ is trivial, and the forcing at stage $\delta_0$ adds a club to $\delta_0$. Thus, the forcing $\mathbb{P}$ up to and including stage $\delta_0$ has cardinality $\delta_0$ and the forcing past stage $\delta_0$ has a dense set which is closed up to the next inaccessible cardinal. Thus, $\mathbb{P}$ has a closure point at $\delta_0$, thus $V \subseteq V[G]$ satisfies the $\delta_0^+$ approximation and cover properties by [Hamkins 11 (Lemma 13)]. Thus, by Lemma 20, the forcing $\mathbb{P}$ does not increase Mitchell rank for $\theta$-supercompactness for cardinals above $\delta_0$ where $\theta$ is also above $\delta_0$. 

Let $\kappa$ be a $\Theta(\kappa)$-supercompact cardinal such that $\Theta''\kappa \subseteq \kappa$, and $\Theta \upharpoonright \kappa$ represents $\theta(\kappa)$, and $F \upharpoonright \kappa$ represents $F(\kappa)$ in $V$. Here we will need that representing functions are still representing functions in the extension. The fact we will use for this is from [Hamkins 11] which says that if $V \subseteq V[G]$ satisfies the hypothesis of the approximation and cover theorem, then every $\Theta(\kappa)$-closed embedding $j:V[G] \to M[j(G)]$ is a lift of an embedding $j \upharpoonright V : V \to M$ that is in $V$ (and $j \upharpoonright V$ is also a $\Theta(\kappa)$-closed embedding). Thus, if $F \upharpoonright \kappa$ is a representing function for $F(\kappa)$ in the ground model, $(j \upharpoonright V)(F \upharpoonright \kappa)(\kappa) = F(\kappa)$, and thus $F \upharpoonright \kappa$ is a representing function for $F(\kappa)$ with respect to such embeddings $j$
 in $V[G]$. 

 The induction hypothesis is that any $\kappa' < \kappa$, where $\Theta''\kappa' \subseteq \kappa'$, and $\Theta \upharpoonright \kappa'$ represents $\Theta(\kappa)$, and $F \upharpoonright \kappa'$ represents $F(\kappa')$ in $V$, has $o_{\Theta(\kappa)-sc}(\kappa')^{V[G]} = \min\{o_{\Theta(\kappa)-sc}(\kappa')^V, F(\kappa')\}$. Pick $j:V \to M$ a $\Theta(\kappa)$-supercompactness embedding with critical point $\kappa$ such that $o_{\Theta(\kappa)-sc}(\kappa)^M = \beta_0 < \min\{o_{\Theta(\kappa)-sc}(\kappa)^V, F(\kappa)\}$. We will lift $j$ through $\mathbb{P}$. The forcing above $\kappa$ does not affect the Mitchell rank for $\Theta(\kappa)$-supercompactness of $\kappa$ since the forcing past stage $\kappa$ has a dense subset which is closed up to the next closure point of $\Theta$ past $\kappa$. Call $\mathbb{P}_{\kappa}*\dot{\mathbb{Q}}$ the forcing up to and including stage $\kappa$.

 First we will lift $j$ through $\mathbb{P}_{\kappa}$. Let $G_{\kappa} \subset \mathbb{P}_{\kappa}$ be $V$-generic and $g \subseteq \mathbb{Q}$ be $V[G_{\kappa}]$-generic. First, we need an $M$-generic filter for $j(\mathbb{P}_{\kappa})$. Since cp$(j) = \kappa$, the forcings $\mathbb{P}_{\kappa}$ and $j(\mathbb{P}_{\kappa})$ agree up to stage $\kappa$. The forcing $\mathbb{Q}$ adds a club $c \subseteq \kappa$ such that $\forall \delta < \kappa, \delta \in c$ implies $o_{\Theta(\kappa)-sc}(\delta)^V < F(\delta)$. Since $o_{\Theta(\kappa)-sc}(\kappa)^V < F(\kappa)$ implies $o_{\Theta(\kappa)-sc}(\kappa)^M < F(\kappa) = j(F)(\kappa)$, and $j(\Theta)(\kappa) = \Theta(\kappa)$, it follows that the $\kappa$th stage of $j(\mathbb{P}_{\kappa})$ is $\mathbb{Q}$. Thus, $j(\mathbb{P}_{\kappa})$ factors as $\mathbb{P}_{\kappa}*\dot{\mathbb{Q}}*\mathbb{P}_{\mbox{tail}}$ where $\mathbb{P}_{\mbox{tail}}$ is the forcing past stage $\kappa$. Since $G_{\kappa}*g$ is $M$-generic for $\mathbb{P}_{\kappa}*\dot{\mathbb{Q}}$, we only need an $M[G_{\kappa}][g]$-generic filter for $\mathbb{P}_{\mbox{tail}}$ which is in $V[G_{\kappa}][g]$. We will diagonalize to obtain such a filter, after checking that we have met the criteria. Since $|\mathbb{P}_{\kappa}|=|\mathbb{Q}| = \kappa$, and $V \models $ GCH, it follows that the forcing $\mathbb{P}_{\kappa}*\dot{\mathbb{Q}}$ has the $\kappa^+$-chain condition. Since $M^{\Theta(\kappa)} \subseteq M$, it follows [Hamkins 8 (Theorem 54)] that $M[G_{\kappa}][g]^{\Theta(\kappa)} \subseteq M[G_{\kappa}][g]$. Since $|\mathbb{P}_{\kappa}| = \kappa$, and has the $\kappa^+$-chain condition, it has at most $\kappa^{<\kappa^+} = \kappa^+$ many maximal antichains. Since $|j(\kappa^+)| \le \kappa^{+^{\Theta(\kappa)^{< \kappa}}} = \Theta(\kappa)^+$, it follows that $\mathbb{P}_{\mbox{tail}}$ has at most $\Theta(\kappa)^+$ many maximal antichains in $M[G_{\kappa}][g]$. Since $\mathbb{P}_{\mbox{tail}}$ has a dense subset which is closed up to the next closure point of $\Theta$ past $\kappa$, it has a dense subset which is $\le \Theta(\kappa)$-closed. Thus, we can enumerate the maximal antichains of $\mathbb{P}_{\mbox{tail}}$, and build a descending sequence of conditions meeting them all through the dense set which has $\le \Theta(\kappa)$-closure, using this closure and the fact that $M[G_{\kappa}][g]$ is closed under $\Theta(\kappa)$-sequences to build through the limit stages. Then the upward closure of this sequence $G_{\mbox{tail}} \subseteq \mathbb{P}_{\mbox{tail}}$ is $M[G_{\kappa}][g]$-generic and is in $V[G_{\kappa}][g]$. Thus, $j(G_{\kappa}) = G_{\kappa}*g*G_{\mbox{tail}}$ is $M$-generic, it is in $V[G_{\kappa}][g]$ and $j''G_{\kappa} \subseteq j(G_{\kappa})$. Thus, we have the partial lift $j: V[G_{\kappa}] \to M[j(G_{\kappa})]$.

 Next we lift though $\mathbb{Q}$ by finding an $M[j(G_{\kappa})]$-generic filter for $j(\mathbb{Q})$. The forcing $j(\mathbb{Q})$ adds a club $C \subseteq j(\kappa)$ such that $\forall \delta < j(\kappa), \delta \in C$ implies $o_{\Theta(\kappa)-sc}(\delta)^{M[j(G_{\kappa})]} <j(F)(\delta)$. Since $j(G_{\kappa}) \in V[G_{\kappa}][g]$ and $M^{\Theta(\kappa)} \subseteq M$, it follows that $M[j(G_{\kappa})]^{\Theta(\kappa)} \subseteq M[j(G_{\kappa})]$ by [Hamkins 8 (Theorem 53)]. Since for every $\beta < \kappa$, there is a dense subset of $\mathbb{Q}$ which is $\le \beta$-closed, it follows that there is a dense subset of $j(\mathbb{Q})$ which is $\le \Theta(\kappa)$-closed. Also, since $|\mathbb{Q}| = \kappa$ and $\mathbb{Q}$ has the $\kappa^+$-chain condition, the forcing $\mathbb{Q}$ has at most $\kappa^+$ many maximal antichains. Therefore $j(\mathbb{Q})$ has at most $|j(\kappa^+)| \le \kappa^{+^{\Theta(\kappa)^{<\kappa}}} = \Theta(\kappa)^+$ many maximal antichains in $M[j(G_{\kappa})]$. Thus we can diagonalize to get a generic filter with a master condition as follows. Let $c = \cup g$ be the new club of $\kappa$ and consider $\overline{c} = c \cup \{ \kappa \}$ which is in $M[j(G_{\kappa})]$. We have $\delta \in c$ implies $o_{\Theta(\kappa)-sc}(\delta)^V < F(\delta)$. Thus, $\delta \in c$ implies $o_{\Theta(\kappa)-sc}(\delta)^M < F(\delta)$. By Lemma 20, it follows that $\delta \in c$ implies $o_{\Theta(\kappa)-sc}(\delta)^{M[j(G_{\kappa})]} \le o_{\Theta(\kappa)-sc}(\delta)^M < F(\delta) = j(F)(\delta)$. Since $o_{\Theta(\kappa)-sc}(\kappa)^{M[j(G)]} \le o_{\Theta(\kappa)-sc}(\kappa)^M = \beta_0 < \min\{o_{\Theta(\kappa)-sc}(\kappa)^V, F(\kappa)\}$, and $j(\Theta)(\kappa) = \Theta(\kappa)$, it follows that $o_{j(\Theta)(\kappa)-sc}(\kappa)^{M[j(G)]} < F(\kappa) = j(F)(\kappa)$. Thus $\overline{c}$ is a closed, bounded subset of $j(\kappa)$ such that every $\delta$ in $\overline{c}$ has $o_{j(\Theta)(\kappa)-sc}(\delta)^{M[j(G_{\kappa})]} < j(F)(\kappa)$. Thus $\overline{c} \in j(\mathbb{Q})$. Thus diagonalize to get an $M[j(G_{\kappa})]$-generic filter $g^* \subseteq j(\mathbb{Q})$ which contains the master condition $\overline{c}$. Then $j''g \subseteq j(g) = g^*$. Therefore we may lift the embedding to $j:V[G_{\kappa}][g] \to M[j(G_{\kappa})][j(g)]$.

 The success of the lifting arguments for $o_{\Theta(\kappa)-sc}(\kappa)^M = \beta_0 < \min\{o_{\Theta(\kappa)-sc}(\kappa)^V, F(\kappa)\}$, where $\beta_0$ was arbitrary, shows that $o_{\Theta(\kappa)-sc}(\kappa)^V[G] \ge \min\{o_{\Theta(\kappa)-sc}(\kappa)^V, F(\kappa)\}$. By Lemma 20, we have $o_{\Theta(\kappa)-sc}(\kappa)^{V[G]} \le o_{\Theta(\kappa)-sc}(\kappa)^V$. All that remains to see is that $o_{\Theta(\kappa)-sc}(\kappa)^{V[G]} \le F(\kappa)$. Consider any $\Theta(\kappa)$-supercompactness embedding $j:V[G] \to M[j(G)]$ with critical point $\kappa$ in $V[G]$, and the new club $c \subseteq \kappa$, where $\forall \delta < \kappa, \delta \in c$ implies $o_{\Theta(\delta)-sc}(\delta)^V < F(\delta)$. This club is in any normal fine measure on $P_{\kappa}(\Theta(\kappa))$. By the induction hypothesis, $\forall \delta < \kappa, \delta \in c$ implies $o_{\Theta(\delta)-sc}(\delta)^{V[G]} < F(\delta)$. Applying $j$ to this statement gives $\forall \delta < j(\kappa), \delta \in j(c)$ implies $o_{j(\Theta)(\kappa)-sc}(\delta)^{M[j(G)]} < j(F)(\delta)$. Since $F \upharpoonright \kappa$ represents $F(\kappa)$, we have $j(F)(\kappa) = F(\kappa)$, and since $\Theta \upharpoonright \kappa$ represents $\Theta(\kappa)$, we have $j(\Theta)(\kappa) = \Theta(\kappa)$. Since $\kappa < j(\kappa)$ and $\kappa \in j(c)$, it follows that $o_{\Theta(\kappa)-sc}(\kappa)^{M[j(G)]} < F(\kappa)$. Thus by Lemma 19, since $j$ was arbitrary, $o_{\Theta(\kappa)-sc}(\kappa)^{V[G]} \le F(\kappa)$. Thus $o_{\Theta(\kappa)-sc}(\kappa)^{V[G]} \le \min\{o_{\Theta(\kappa)-sc}(\kappa)^V, F(\kappa)\}.$ Thus for any $\kappa$ which is $\Theta(\kappa)$-supercompact for which $\Theta''\kappa \subseteq \kappa$, and $F \upharpoonright \kappa$ represents $F(\kappa)$, and $\Theta \upharpoonright \kappa$ represents $\Theta(\kappa)$ in $V$ has $o_{\Theta(\kappa)-sc}(\kappa)^{V[G]} = \min\{o_{\Theta(\kappa)-sc}(\kappa)^V, F(\kappa)\}$.
 \end{proof}

\end{document}
